\subsection{The exact higher central lifting image, including convergence and Frobenius corrections}

\begin{statement}{E.44}{Proved}{A global higher lift is equivalent to integral restricted tame corrections}
Let $W$ be any of the three actual cubic quotients in
(E41.1), and retain $W_r$ and $\mathcal A_r(W)$ from
(E41.2). Fix $r\ge3$, put $A=\mathcal A_r(W)$, and
let $a:G_{\mathbb Q}\to W_{r-1}$ be an actual continuous
cocycle for the based arithmetic action. It may be
ramified outside three. Fix the continuous Hall set
section $h:W_{r-1}\to W_r$ with $h(1)=1$ from E.41,
and write $\widehat a=h\circ a$.

For every prime $\mathfrak q\nmid3$ of
$K=\mathbb Q(\zeta)$ fix one actual prolongation,
a tame pro-three inertia generator $t_{\mathfrak q}$,
and an arithmetic Frobenius $\phi_{\mathfrak q}$.
Write $q=N\mathfrak q$; the tame relation is
$\phi_{\mathfrak q}t_{\mathfrak q}\phi_{\mathfrak q}^{-1}
=t_{\mathfrak q}^q$ after killing wild inertia.
Set
\[
 \begin{split}
 n_{\mathfrak q}&=\widehat a(t_{\mathfrak q}),\qquad
 x_{\mathfrak q}=\widehat a(\phi_{\mathfrak q}),\\
 \epsilon_{\mathfrak q}
  &=x_{\mathfrak q}\,
      \rho_r(\phi_{\mathfrak q})(n_{\mathfrak q})\,
      x_{\mathfrak q}^{-1}n_{\mathfrak q}^{-q}\ \in A,\\
 z_{\mathfrak q}
  &=-(\Phi_{\mathfrak q}-qI)^{-1}\epsilon_{\mathfrak q}
                                  \ \in A\otimes\mathbb Q_3.
 \end{split}
 \tag{E44.1}
\]
Here $\rho_r$ is the genuine action on $W_r$ and
$\Phi_{\mathfrak q}$ is its induced linear action
on $A$. The central element $\epsilon_{\mathfrak q}$
is written additively in the last line. The inverse
exists for every prime in (E44.1).

The geometric-conjugacy fiber of the actual map
\[
 \operatorname{Sect}(W_r\rtimes G_{\mathbb Q})/W_r
       \longrightarrow
 \operatorname{Sect}(W_{r-1}\rtimes G_{\mathbb Q})/W_{r-1}
 \tag{E44.2}
\]
over $[a]$ is nonempty exactly for the following
explicit condition:
\[
 z_{\mathfrak q}\in A\quad(\mathfrak q\nmid3),
 \qquad
 \text{for every }m,\quad
 z_{\mathfrak q}\in3^mA
       \text{ outside finitely many }\mathfrak q.
 \tag{E44.3}
\]
Every nonempty fiber is a torsor under the evaluated
finite-matrix group $\mathcal C_r(W)$ in (E41.3).
Thus the coefficient lattice and all parameters
within a nonempty fiber are already explicit.

More precisely, write
\[
 \mathcal R(A)=\{(v_{\mathfrak q})\in\prod_{\mathfrak q\nmid3}A:
                  v_{\mathfrak q}\longrightarrow0
                         \text{ in the three-adic topology}\},
 \quad
 d(v)_{\mathfrak q}=(q^{-1}\Phi_{\mathfrak q}-1)v_{\mathfrak q}.
\]
The exact continuous localization map sends the
actual obstruction to
\[
 H^2_{\rm cts}(K,A)
      \xrightarrow{\ \sim\ }\mathcal R(A)/d\mathcal R(A),
 \qquad
 \operatorname{ob}_r(a|_{G_K})
      \longmapsto[(q^{-1}\epsilon_{\mathfrak q})_{\mathfrak q}].
 \tag{E44.4}
\]
The denominator is the image, with no closure
added. The restricted sequence condition in
(E44.3) consequently retains the possible
inverse-limit obstruction.

One local target in this criterion is completely
evaluated in the next intrinsic degree. At the
$K$-prime $(2)$, with $q=4$, the common degree-three
module of all three quotients has
\[
 \operatorname{coker}(\Phi_{(2)}-4I)
    =(\mathbb Z/3)^{70}\oplus(\mathbb Z/9)^{24}
      \oplus(\mathbb Z/27)^{76}\oplus(\mathbb Z/81)^{28}.
 \tag{E44.5}
\]
The group has order $3^{458}$. At the rational
prime two the corresponding actual target is
\[
 H^2_{\rm cts}(\mathbb Q_2,\operatorname{Lie}_3V)
   =\operatorname{coker}(\Phi_{2,\mathbb Q}-2I)
   =(\mathbb Z/3)^{36}\oplus(\mathbb Z/9)^{12}
     \oplus(\mathbb Z/27)^{38}\oplus(\mathbb Z/81)^{14},
 \tag{E44.5a}
\]
of order $3^{230}$. Their Smith certificates
and matrices are included in
\texttt{e\_higher\_images.py}, its output and JSON file.

The error in (E44.1) uses the actual higher
Frobenius action. With the splitting in E.39,
its lower Frobenius value is exactly
\[
 a(\phi_{\mathfrak q})
      =\lambda(h_{\mathfrak q})u(\sigma_{\mathfrak q}).
 \tag{E44.6}
\]
Formula (E44.4) proves the higher lifting image
for these supplied arithmetic data; evaluation
of all the higher words in (E44.1) and (E44.6)
for all ramified inputs remains the further
step in the full metabelian and marked section
classification. The rank-ten linear matrices
determine $\Phi_{\mathfrak q}$, while
$\rho_r(\phi_{\mathfrak q})$ is the genuine
action on the specified group quotient.

\emph{Manuscript: \S51.16, folios 733v--736v,
scans 172--178, (371)--(379), Part E(c),(d).
The comparison uses the intrinsic quotients
and the exact maps in (E41.23); it applies
above ramified as well as unramified lower
sections.}
\end{statement}
\begin{proof}
\emph{The actual central extension and its local error.}
The kernel of $W_r\to W_{r-1}$ is the finite
free lattice $A$ proved in E.41. It is central
in $W_r$, so pullback along any lower section
has its original arithmetic action on $A$.
The continuous set lift $\widehat a$ defines
the central continuous two-cocycle
\[
 \omega(g,h)=\widehat a(g)\,
       g(\widehat a(h))\,\widehat a(gh)^{-1}\in A.
 \tag{E44.7}
\]
Associativity gives its cocycle identity,
and changing the set lift changes it by
a continuous coboundary. A continuous
$A$-valued correcting cochain $b$ with
$\delta b=-\omega$ gives the actual lifted
cocycle $\widetilde a(g)=b(g)\widehat a(g)$.
Conversely every lift has this form. Thus
the vanishing of the continuous class of
(E44.7) is exactly the lifting assertion.

At $\mathfrak q\nmid3$ the based arithmetic
action on every $W_r$ is unramified. Wild
inertia has trivial cocycle image: it is
pro-$p$ for $p\ne3$ and the geometric
target is pro-three. On tame inertia the
lower cocycle is therefore a homomorphism,
and
$a(t_{\mathfrak q}^s)=a(t_{\mathfrak q})^s$
for $s\in\mathbb Z_3$. The corresponding
lifted inertia homomorphism is
\[
 t_{\mathfrak q}^s\longmapsto n_{\mathfrak q}^s.
 \tag{E44.8}
\]
Its difference from
$\widehat a(t_{\mathfrak q}^s)$ is an
explicit local null-homotopy of (E44.7).
The lower section satisfies the tame
relation, so the error in (E44.1) belongs
to the central kernel $A$.

Replacing $n_{\mathfrak q}$ by
$z n_{\mathfrak q}$ with $z\in A$ changes
this error by exactly
\[
 \epsilon_{\mathfrak q}\longmapsto
    \epsilon_{\mathfrak q}+(\Phi_{\mathfrak q}-qI)z.
 \tag{E44.9}
\]
This follows by moving the central factors
through the displayed word: Frobenius
contributes $\Phi_{\mathfrak q}z$ and the
$q$th power contributes $qz$.
Changing $x_{\mathfrak q}$ by a central
element has no effect at all. Thus (E44.1)
is the unique rational central correction
of the chosen inertia lift.

The weights on $A\otimes\mathbb Q_3$ are
among the degree-$r$ tensor weights of
$V$, namely Frobenius absolute values
$q^{j/2}$ with $r\le j\le2r$.
This follows from each actual graded
quotient in (E41.5). For $r\ge3$ none is
$q$, so $\Phi_{\mathfrak q}-qI$ is
invertible over $\mathbb Q_3$. There is
no choice of a rational kernel coordinate
in (E44.1).

\emph{Localization at finite coefficients, followed by the derived limit.}
Put $U=\operatorname{Spec}\mathcal O_K[1/3]$
and $A_m=A/3^mA$. The action on this finite
module over $G_{K,\{3\}}$ has three-group
image by E.34. A finite three-group module
has a filtration by trivial
$\mathbb F_3$ modules. The vanishing
$H^2(G_{K,\{3\}},\mathbb F_3)=0$ proved
in E.34 consequently gives
\[
 H^2(U,A_m)=0,\qquad H^3(U,A_m)=0,
             \qquad H^1(U,A_m)\text{ finite}.
 \tag{E44.10}
\]
Here the comparison is
\cite[II.2.9]{E_MilneADT}; the degree-three
vanishing and finiteness are
\cite[I.4.10(c), I.4.15]{E_MilneADT}.
The field $K$ has no real places and
three is inverted on $U$, as these
statements require.

For a finite set $J$ of primes of $U$,
the actual localization triangle is
\[
 R\Gamma(U,A_m)\longrightarrow
 R\Gamma(U-J,A_m)\longrightarrow
 \bigoplus_{\mathfrak q\in J}
       R\Gamma(k_{\mathfrak q},A_m(-1))[-1]
 \longrightarrow R\Gamma(U,A_m)[1].
 \tag{E44.11}
\]
It follows from
\cite[II.0.1(f), II.1.7(b)]{E_MilneADT}.
The purity term can also be computed
directly over the strict Henselization:
wild inertia contributes no three-primary
cohomology, tame inertia is
$\mathbb Z_3(1)$, and its degree-one
cohomology is $A_m(-1)$.
Taking the filtered direct limit over
finite $J$ gives the generic field
$G_K$ term. The support term remains
a direct sum at each fixed coefficient
level.

Now take $R\varprojlim_m$. Its comparison
with continuous cochains and its derived
limit sequence are
\cite[Theorem 2.2 and Proposition 1.6]{E_JannsenContinuous}.
Equation (E44.10) and the finiteness of
$H^1(U,A_m)$ give zero in degrees two
and three for the unramified inverse-limit
complex. For the residue-field term use
the compatible two-term complex
\[
 [\,A_m\xrightarrow{q^{-1}\Phi_{\mathfrak q}-1}A_m\,]
                \quad\text{in degrees }0,1.
 \tag{E44.12}
\]
The direct-sum transitions in both terms
are surjective. Their derived inverse
limit is therefore the literal complex
\[
 [\,\mathcal R(A)\xrightarrow{d}\mathcal R(A)\,].
 \tag{E44.13}
\]
The resulting degree-two isomorphism is
the first part of (E44.4). In particular
the image in its denominator is retained
as an image; this calculation has not
replaced inverse limits of complexes by
inverse limits of cohomology groups.

\emph{The chain-level representative and uniform continuity.}
The support map in (E44.11) records the
failure of an inertia splitting to be
Frobenius invariant. Use exactly the
splitting (E44.8). Transporting it by
$(x_{\mathfrak q},\phi_{\mathfrak q})$
and evaluating at $t_{\mathfrak q}$
gives the error at $t_{\mathfrak q}^q$
pulled back by its $q$th-root map.
In formulas,
\[
 x_{\mathfrak q}\phi_{\mathfrak q}(n_{\mathfrak q})
          x_{\mathfrak q}^{-1}
       =\epsilon_{\mathfrak q}n_{\mathfrak q}^q
 \quad\Longrightarrow\quad
 x_{\mathfrak q}\phi_{\mathfrak q}(n_{\mathfrak q}^{q^{-1}})
          x_{\mathfrak q}^{-1}n_{\mathfrak q}^{-1}
       =q^{-1}\epsilon_{\mathfrak q}.
 \tag{E44.14}
\]
This is the twist in
$\operatorname{Hom}(\mathbb Z_3(1),A)=A(-1)$.
Equation (E44.9) divided by $q$ is
exactly the differential in (E44.12).
Thus the representative is
$D_{\mathfrak q}=q^{-1}\epsilon_{\mathfrak q}$
at the cochain level, with the sign in
(E44.1) forced by (E44.9).

We verify that these local null-homotopies
are admissible in the direct-sum complexes,
including their inverse limit. For every
finite characteristic quotient of
$W_{r-1}$, the map
$g\mapsto(a(g),\rho(g))$ modulo that
quotient has finite image. Its kernel
defines a finite extension of $K$,
ramified at finitely many primes.
Hence $a(t)$ approaches $1$ uniformly
on every inertia group outside a finite
set, including all prolongations of those
primes. Continuity of $h$ at $1$ gives
the same assertion for $h(a(t))$ in
every finite quotient of $W_r$.
The map $(n,s)\mapsto n^s$ is continuous
on the compact space $W_r\times\mathbb Z_3$;
indeed a closed subgroup containing $n$
contains every $n^s$.
It follows that the difference between
$n_{\mathfrak q}^s$ and
$h(a(t_{\mathfrak q}^s))$ is zero
modulo $3^mA$, uniformly in $s$, outside
finitely many base primes.

The conjugation and arithmetic action
are continuous and preserve sufficiently
small characteristic open subgroups.
The same argument proves
$\epsilon_{\mathfrak q}\in3^mA$
outside finitely many primes. On the
prolongations of any fixed prime, these
finite-coefficient cochains depend
continuously, and hence locally
constantly, on the conjugating Galois
element. In (E44.4) we use the one
fixed actual prolongation per base
prime specified in the statement.
These facts justify the direct-sum
representative before taking the
derived limit and prove
$(D_{\mathfrak q})\in\mathcal R(A)$.
Equations (E44.11)--(E44.14) now give
the second part of (E44.4), including
its actual representative in the
restricted quotient.

The component maps $d_{\mathfrak q}$
are rationally injective. Consequently
$D$ is in $d\mathcal R(A)$ precisely
when its unique negative inverse is
the integral restricted sequence in
(E44.3). This proves the criterion
over $K$. Restriction from
$H^2_{\rm cts}(\mathbb Q,A)$ to
$H^2_{\rm cts}(K,A)$ is injective:
the composite with corestriction is
multiplication by two, which is an
automorphism of this three-primary
cohomology module. Thus the criterion
also proves the asserted global lift
over $\mathbb Q$.

\emph{Actual gluing and independence of choices.}
The vanishing of the continuous class
in (E44.7) supplies a continuous
global correcting cochain, hence an
actual lift. The gluing can also be
performed at finite coefficients.
The reductions of the corrected
inertia lifts
$z_{\mathfrak q}n_{\mathfrak q}$
are compatible and have finite support
in their central corrections at each
level. Finite-coefficient localization
first gives a global lift modulo
$3^mA$. Its difference from the
prescribed inertia lifts is a finite
family in
$\bigoplus_{\mathfrak q}H^0(k_{\mathfrak q},A_m(-1))$.
The preceding degree of (E44.11) is
surjective onto this family because
$H^2(U,A_m)=0$. Subtract the resulting
global one-cocycle to align every
prescribed inertia lift.

Two such aligned global correcting
cochains differ by a cocycle zero on
all these inertia groups, hence by
an unramified cocycle. The group of
such finite-coefficient cocycles is
finite: its cohomology is finite by
(E44.10), and its coboundaries are
images of the finite module $A_m$.
The nonempty finite sets of aligned
cochains therefore have a compatible
inverse limit. Each component is
continuous, so the resulting
$A$-valued cochain is continuous.
This construction also shows that
the continuous inertia-bundle
equivariance of E.39 is fulfilled by
the global lift; no independent
choice of incoherent inertia values
has been introduced.

Changing the Hall lift changes
$n_{\mathfrak q}$ by an integral
restricted central sequence, by
the uniform continuity just proved.
Equation (E44.9) changes $D$ by its
image under $d$, preserving (E44.3).
Changing Frobenius by an inertia
power replaces its geometric lift
by $n_{\mathfrak q}^s x_{\mathfrak q}$
up to a central factor; this leaves
$\epsilon_{\mathfrak q}$ unchanged.
A change of tame generator first transports
the inertia lift to $n_{\mathfrak q}^u$,
$u\in\mathbb Z_3^*$, and rescales
its additive coordinate by that unit
in $\mathbb Z_3$, as does the
identification of the corresponding
Tate line. Such rescalings preserve
integrality and the restricted
condition. Returning from this
transported lift to
$h(a(t_{\mathfrak q}^u))$ changes
the representative by a restricted
central boundary, by the same
uniform continuity argument.
Changing prolongations similarly
gives the corresponding Galois-module
coordinate isomorphisms followed
by a restricted central boundary
when returning to the Hall lift.
Their integral automorphisms preserve
every $3^mA$. Finally, geometric
conjugation of $a$ can be lifted
to $W_r$; it conjugates the error
word and acts trivially on its
central value. Returning to Hall
lifts again adds a restricted
boundary. Thus the criterion is
well defined on the fiber (E44.2).

The differences between two global
lifts are $A$-valued one-cocycles.
At every good prime their inertia
value $v$ satisfies $\Phi v=qv$.
The same weight calculation forces
$v=0$, so every such cocycle is
unramified outside three. The
lowest nonzero graded coordinate
argument of E.41 excludes every
section-twisted lower stabilizer,
including for a ramified section.
The full geometric-conjugacy
fiber is therefore the torsor
$\mathcal C_r(W)$, with no further
quotient. This supplies the claimed
parameters whenever (E44.3) holds.

\emph{An evaluated degree-three local obstruction module.}
At the $K$-prime $(2)$ the actual
elliptic scalar is $-2$ and the
toric scalar is four. In degree
three the split-lattice eigenvalues
and multiplicities are
\[
\begin{array}{c|rrrr}
 \text{eigenvalue}&64&-32&16&-8\\\hline
 \text{multiplicity}&168&128&32&2.
\end{array}
 \tag{E44.15}
\]
They come from zero, one, two and
three elliptic letters, respectively.
The actual rank-$330$ integral
matrix is constructed from the
matrix $A_2$ of E.37 in Hall
coordinates. The split-lattice
maps satisfy $B_3C_3=27I$, and
$9/(\lambda-4)$ is three-integral
for each eigenvalue $\lambda$ in
(E44.15). Hence
\[
 3^5(\Phi_{(2)}-4I)^{-1}A\subset A.
 \tag{E44.16}
\]
The local cokernel is killed by
$3^5$ and has full rational rank
zero. The saved modular row and
column transformations in
\texttt{e\_higher\_images.py}
verify modulo $729$ the exact
Smith valuation counts
\[
 0^{\times132},\quad1^{\times70},\quad
 2^{\times24},\quad3^{\times76},\quad4^{\times28}.
 \tag{E44.17}
\]
All transformations are invertible
modulo three. The proven bound
(E44.16) makes this finite
certificate conclusive. As an
independent check, its valuation
sum is
$70+48+228+112=458$, which equals
the valuation of the determinant
from (E44.15), namely
$168+2\cdot128+32+2=458$.
The three actual group quotients
have identical degree-three
kernels by E.40--E.41, so this
proves (E44.5) for all of them.

We also compute the rational-prime
target, retaining the integral
arithmetic comparison. Let
$\mathcal O=\mathbb Z_3[\zeta]$.
The rational Frobenius at two
acts on the rank-one
$\mathcal O$-module $T_3E$ as
$u c$ for $u\in1+3\mathcal O$:
all $E[3]$ is defined over $K$,
and its nontrivial residual
Galois element is conjugation.
The determinant is two by the
Weil pairing, while $\det(c)=-1$,
so $u\bar u=-2$. The point count
on $Y^2+Y=X^3-7$ over
$\mathbb F_2$ gives three points
and trace zero, in agreement
with its square $-2$ over $K$.

Let $s=\sqrt{-2}\in1+3\mathbb Z_3$
and set $v=(s/u)^{1/2}$ in
$1+3\mathcal O$. Unique square
roots exist there; $v\bar v=1$
and $v u/\bar v=s$. The full
CM image proved in E.37 realizes
$v$ by an actual $G_K$ action
whose toric action is the
identity. Its entire integral
$V$-matrix is
\[
 \begin{pmatrix}
 v&0\\ U(v-I)/3&I_8
 \end{pmatrix}.
 \tag{E44.18}
\]
Thus it conjugates the actual
rational Frobenius representation
to the matrix with elliptic
block $sC_E$, toric block $2P_D$,
and lower block
$(UsC_E-2P_DU)/3$. Here
$P_D=-C_T$ is the cusp
permutation, and $U,C_E,C_T$
are exactly the E.37 matrices.
The comparison (E44.18) acts
on the whole generalized-Jacobian
lattice and is integral.

The square root modulo $2187$
determines this divided block
modulo $729$. The program
constructs its degree-three
matrix, verifies that its square
is the $K$-Frobenius matrix
modulo $729$, and gives Smith
valuations
\[
 0^{\times230},\quad1^{\times36},\quad
 2^{\times12},\quad3^{\times38},\quad4^{\times14}.
 \tag{E44.19}
\]
The exact exponent bound follows
from (E44.16) and
\[
 (\Phi_{2,\mathbb Q}-2I)^{-1}
   =(\Phi_{2,\mathbb Q}+2I)
                     (\Phi_{(2)}-4I)^{-1}.
\]
The matrix on the left is
therefore determined for Smith
purposes at the computed
precision. An independent
determinant check uses the
split degree-one eigenvalues
$s,-s,2^{\times5},(-2)^{\times3}$
and all degree-three Hall
monomials; its valuation is
$36+24+114+56=230$.
The local support sequence has
zero degree-two and degree-three
unramified residue terms, so its
degree-two field cohomology is
precisely the displayed cokernel.
This proves (E44.5a).

The inverse corrections included
in the code are lattice checks;
the theorem does not assert that
every arbitrary lattice error is
realized by a global lower section.

Lastly (E44.6) is the exact
decomposition (E39.5) applied
to the actual local Frobenius.
Substituting it into (E44.1)
retains its correction word and
the genuine higher arithmetic
action. E.38 and E.43 evaluate
particular ramified inputs;
(E44.1)--(E44.4) give the proved
general lifting test with its
remaining arithmetic data
explicitly specified.
\end{proof}
